\title{DeepSeekMath: Pushing the Limits of Mathematical Reasoning in Open Language Models}

\begin{abstract}
Mathematical reasoning remains a particularly demanding domain for language models, owing to its inherent structure and intricacy. This study presents DeepSeekMath~7B, a model obtained by further pre-training DeepSeek-Coder-Base-v1.5 7B using 120B math-centric tokens collected from Common Crawl, augmented by natural language and code corpora. DeepSeekMath~7B attains a notable 51.7\% accuracy on the MATH competition benchmark without leveraging any auxiliary toolkits or ensemble voting strategies, nearly matching the results achieved by Gemini-Ultra and GPT-4. Furthermore, when employing self-consistency with 64 samples, DeepSeekMath~7B reaches a 60.9\% score on the MATH benchmark. The advancement in DeepSeekMath~’s mathematical reasoning is attributed to two primary components: (1) capitalizing on the abundance of open web data through a carefully designed data filtering and curation system; and (2) implementing Group Relative Policy Optimization (GRPO)—an adaptation of Proximal Policy Optimization (PPO)—to foster enhanced reasoning skills while also improving the efficiency of PPO in terms of memory consumption.
\end{abstract}

\section{Introduction}

Large language models (LLMs) have fundamentally transformed mathematical reasoning within artificial intelligence, driving major progress on both quantitative reasoning benchmarks \citep{MATH} and geometry-focused benchmarks \citep{trinh2024solving}. In addition, these systems have become valuable partners in tackling sophisticated mathematical challenges \citep{tao}. Nevertheless, top-performing solutions such as GPT-4 \citep{gpt4} and Gemini-Ultra \citep{gemini} remain proprietary, and open-access models to date significantly lag behind in mathematical problem-solving performance.

This paper presents DeepSeekMath, a specialized language model in the mathematical domain that markedly surpasses existing open-source models and achieves results approaching those of GPT-4 on standard academic assessments. The foundation for this advancement is the construction of the DeepSeekMath Corpus, an extensive and high-caliber pre-training dataset containing 120B mathematical tokens. The data collection process begins by extracting content from Common Crawl (CC) through the deployment of a classifier built on fastText \citep{joulin2016fasttext}. In the first phase, this classifier is trained using positive samples from OpenWebMath \citep{openwebmath} and a broad array of unrelated web pages as negative samples. Leveraging the classifier, further positive examples are identified in CC; these candidates undergo an additional round of refinement via human annotation. The improved dataset is subsequently used to retrain the classifier for enhanced detection accuracy. Our experiments demonstrate that the resulting corpus maintains high quality: the base model DeepSeekMath-Base 7B attains 64.2\% accuracy on GSM8K \citep{gsm8k} and 36.2\% on the MATH competition benchmark \citep{MATH}, thereby outperforming Minerva 540B \citep{minerva}. Furthermore, since the DeepSeekMath Corpus incorporates multiple languages, we observe performance gains on Chinese mathematics evaluations \citep{wei2023cmath,agieval}. We contend that our methodology in compiling mathematical datasets lays an initial foundation for subsequent investigations in the field, with considerable potential for further progress.

DeepSeekMath-Base is built upon DeepSeek-Coder-Base-v1.5 7B \citep{deepseek-coder}, as we found initializing from a code-centric pretrained model outperforms starting with a general-purpose LLM. Notably, we observe that mathematical pre-training also confers performance gains on benchmarks such as MMLU \citep{mmlu} and BBH \citep{bbh}, suggesting the training not only boosts mathematical proficiency but also strengthens the model's general reasoning skills.

Following pre-training, we conduct mathematical instruction tuning for DeepSeekMath-Base, incorporating data representing chain-of-thought reasoning \citep{cot}, program-of-thought approaches \citep{pot,pal}, and tool-augmented reasoning \citep{tora}. The resulting model, DeepSeekMath-Instruct 7B, surpasses all other 7B-scale models and achieves results comparable to instruction-tuned open-source models at the 70B scale.

Additionally, we propose Group Relative Policy Optimization (GRPO), a novel reinforcement learning (RL) method derived from Proximal Policy Optimization (PPO) \citep{schulman2017proximal}. Unlike PPO, GRPO eliminates the need for a critic model, instead leveraging group-based score baselines, which drastically curtails resource requirements during training. By applying GRPO solely to a portion of English instruction tuning data, we realize marked improvements over the already competitive DeepSeekMath-Instruct. This is evident across both in-domain datasets (GSM8K: 82.9\% $\rightarrow$ 88.2\%, MATH: 46.8\% $\rightarrow$ 51.7\%) and out-of-domain tasks such as CMATH (84.6\% $\rightarrow$ 88.8\%) in the reinforcement learning stage. We further develop a unified framework that integrates several approaches—including Rejection Sampling Fine-Tuning (RFT) \citep{yuan2023scaling}, Direct Preference Optimization (DPO) \citep{dpo}, PPO, and GRPO—demonstrating that all these strategies can be framed as direct or streamlined RL methodologies. Through a series of comprehensive experiments—including comparisons of online vs. offline training, outcome-based vs. process-based supervision, and single-turn vs. iterative RL—we systematically analyze the key factors underpinning this unified framework. Finally, we elucidate the reasons behind the performance improvements imparted by RL to instruction-tuned models and highlight promising avenues for advancing RL methods within this paradigm.

\subsection{Contributions}
This work makes advances in large-scale mathematical pre-training and delves into the investigation of reinforcement learning approaches.

\textbf{Math Pre-Training at Scale}

\begin{itemize}[topsep=0pt]
    \item Our study demonstrates that the Common Crawl dataset, which is publicly available, contains significant mathematical resources. Leveraging a carefully engineered data filtering process, we assemble the DeepSeekMath Corpus—a rigorously curated collection encompassing 120B tokens extracted from web sources with mathematical content. This dataset substantially surpasses previous efforts, being nearly 7 times greater than Minerva’s math web collection \citep{minerva} and 9 times larger than the recently introduced OpenWebMath dataset \citep{openwebmath}.

    \item DeepSeekMath-Base 7B, our base pre-trained model, matches the performance levels of Minerva 540B \citep{minerva}. This outcome suggests that scaling model parameters is not the sole driver of mathematical reasoning strength; high-quality, targeted pre-training allows even more compact models to excel.

    \item Our math-focused pre-training trials reveal additional insights. Initial training on code prior to math-related content notably enhances the model’s capacity to address mathematical problems, with or without external tools. These results contribute to the ongoing debate regarding whether code pre-training bolsters reasoning performance. For mathematics, our results affirm that it does.

    \item While incorporating arXiv paper content is a widespread practice in mathematics model development, our experiments show it produces no appreciable gains across the suite of mathematical benchmarks evaluated herein.
\end{itemize}

\textbf{Exploration and Analysis of Reinforcement Learning}

\begin{itemize}[topsep=0pt]
    \item We present Group Relative Policy Optimization (GRPO), a reinforcement learning algorithm designed for both efficiency and efficacy. Unlike standard approaches such as Proximal Policy Optimization (PPO), GRPO eliminates the need for a critic network by leveraging group scores to estimate the baseline, which leads to a substantial reduction in the computational overhead during training.
    \item Our results indicate that applying GRPO to instruction tuning substantially improves the performance of our model, DeepSeekMath-Instruct, relying exclusively on instruction-tuning data. Additionally, the reinforcement learning procedure contributes to increased robustness and improved out-of-domain results.
    \item We propose a unified perspective that encompasses a range of methods, including RFT, DPO, PPO, and GRPO. Through an extensive suite of experiments—covering online versus offline training, outcome versus process supervision, and comparisons between single-turn and iterative reinforcement learning—we analyze the critical factors underlying this unified framework.
    \item Building on this framework, we delve into the mechanisms driving the efficacy of reinforcement learning and outline several promising avenues for further advancing reinforcement learning for LLMs.
\end{itemize}

\subsection{Summary of Evaluations and Metrics}

\begin{itemize}[topsep=0pt]
    \item \textbf{English and Chinese Mathematical Reasoning}: Our evaluation comprehensively examines the mathematical problem-solving ability of our models across a spectrum of English and Chinese datasets, spanning tasks from elementary to collegiate complexity. The English evaluation suite comprises GSM8K \citep{gsm8k}, MATH \citep{MATH}, SAT \citep{llemma}, OCW Courses \citep{minerva}, and MMLU-STEM \citep{mmlu}, while for Chinese benchmarks we utilize MGSM-zh \citep{mgsm}, CMATH \citep{wei2023cmath}, Gaokao-MathCloze \citep{agieval}, and Gaokao-MathQA \citep{agieval}. Performance is assessed by the capability to generate self-contained, tool-free solutions as well as the proficiency in solving problems programmatically with Python.

    On the English datasets, DeepSeekMath-Base exhibits performance comparable to the proprietary Minerva 540B \citep{minerva}, outperforming all open-source base models—including Mistral 7B \citep{mistral} and Llemma-34B \citep{llemma}—by a considerable margin, regardless of prior math-specific pre-training. DeepSeekMath-Base's superiority is even more pronounced in the Chinese context, attributable in part to our divergence from the English-only pre-training approach adopted by earlier studies \citep{minerva,llemma}, as we also incorporate diverse, high-quality non-English sources. After incorporating mathematical instruction tuning and reinforcement learning, our DeepSeekMath-Instruct and DeepSeekMath-RL models achieve unprecedented open-source accuracy, surpassing the 50\% threshold on the challenging MATH benchmark.
\end{itemize}

\item \textbf{Formal Mathematics}: DeepSeekMath-Base is assessed on its ability to convert informal mathematical statements into formal proofs, utilizing the evaluation framework established in \citep{dsp_proof} and the miniF2F dataset \citep{minif2f}, with Isabelle \citep{isabelle} serving as the proof assistant. The model achieves notable performance in few-shot settings for automatic formalization.

\item \textbf{Natural Language Understanding, Reasoning, and Code}: In order to thoroughly gauge the model's broad linguistic, reasoning, and coding proficiency, DeepSeekMath-Base is benchmarked on the Massive Multitask Language Understanding (MMLU) suite \citep{mmlu}, which comprises 57 multiple-choice problems spanning a wide array of disciplines; BIG-Bench Hard (BBH) \citep{bbh}, which features 23 tasks that predominantly demand complex, multi-step reasoning; along with HumanEval \citep{codex} and MBPP \citep{mbpp}, two standard metrics for evaluating code generation models. Results indicate that mathematical pre-training enhances the model's capabilities not only in understanding language but also in advanced reasoning tasks.
\end{itemize}

\section{Math Pre-Training}

\subsection{Data Collection and Decontamination}
This section details the procedure employed to assemble the DeepSeekMath Corpus, which sources data from Common Crawl. Figure \ref{fig:data_collect} illustrates our iterative data acquisition pipeline, which effectively accumulates a large mathematical corpus starting from a carefully curated seed dataset, such as a limited but high-quality set of mathematical texts. This pipeline is adaptable and can be employed in other domains, for instance, for collecting programming data.

The process commences with OpenWebMath \citep{openwebmath}, a high-quality corpus of mathematical web documents, selected as the seed. With this resource, we train a fastText classifier \citep{joulin2016fasttext} to identify additional mathematical pages that are stylistically and topically similar to OpenWebMath. To do so, we sample 500,000 entries from the seed corpus as positive examples and pair these with 500,000 randomly selected Common Crawl web pages as negatives. Training is performed using a publicly available fastText implementation\footnote{\url{https://fasttext.cc}}, specifying a 256-dimensional vector space, a learning rate of 0.1, a maximum n-gram length of 3, a minimum word frequency of 3, and three training epochs. To narrow down the massive Common Crawl dataset, both URL-based deduplication and near-duplicate elimination are carried out, yielding a subset of 40 billion unique HTML pages. Utilizing the trained fastText model, mathematical pages are retrieved from this deduplicated set. To ensure content quality, the retrieved documents are ranked by the fastText score, retaining only the highest scoring subset. The quantity of preserved data is subsequently evaluated via pre-training on samples of the top 40B, 80B, 120B, and 160B tokens; for the initial cycle, the decision is made to retain only the leading 40B tokens.

Following the initial phase of data gathering, a considerable portion of mathematical web pages remains uncollected, primarily due to the limited diversity of positive samples used in training the fastText model. To address this, we seek to expand the set of initial sources by identifying additional mathematics-focused web domains to enhance the seed corpus and thereby improve the model's recall. The procedure begins with partitioning the entire Common Crawl dataset into non-overlapping domains, each comprising web pages that share the same base URL. For each domain, we compute the proportion of pages retrieved in the initial collection cycle. If more than 10\% of a domain’s pages were obtained, we categorize it as math-related (e.g., \url{mathoverflow.net}).

Next, we manually label URLs within these domains that are associated with mathematical topics (for example, \url{mathoverflow.net/questions}). Any mathematical pages linked to these labeled URLs that were not previously collected are then added to the seed dataset. By incorporating these additional positive examples, the subsequent training of fastText yields a more robust model, better able to capture mathematical web pages in later rounds. Through four cycles of this collection process, we amass a total of 35.5 million mathematical web pages, containing 120 billion tokens. Observing that 98\% of content gathered in the fourth cycle was already included by the third, we conclude the data collection process at this stage.

In order to prevent overlap with established benchmarks, we adopt the filtering protocol outlined in \cite{deepseek-coder}, removing pages that include questions or answers present in well-known English mathematical benchmarks (GSM8K \citep{gsm8k}, MATH \citep{MATH}) and Chinese datasets (CMATH \citep{wei2023cmath}, AGIEval \citep{agieval}). Specifically, we eliminate from the training set any text segment where a contiguous 10-token substring is an exact match to any excerpt from these evaluation datasets. For benchmark excerpts shorter than 10 tokens but with a minimum length of 3 tokens, exact string matching is also employed to discard contaminated documents.

\subsection{Assessing the DeepSeekMath~Corpus Quality}

We conduct pre-training experiments to compare the DeepSeekMath~Corpus against recently introduced mathematics-focused corpora:

\begin{itemize}[topsep=0pt]
    \item \textbf{MathPile} \citep{mathpile}: a corpus aggregating 8.9B tokens from a variety of sources—including textbooks, Wikipedia, ProofWiki, CommonCrawl, StackExchange, and arXiv—with arXiv making up the majority share (over 85\%);
    \item \textbf{OpenWebMath} \citep{openwebmath}: CommonCrawl-derived data specifically filtered for mathematical content, comprising 13.6B tokens;
    \item \textbf{Proof-Pile-2} \citep{llemma}: this corpus integrates OpenWebMath, AlgebraicStack (10.3B tokens of mathematical code), and arXiv (28.0B tokens); in our experiments with Proof-Pile-2, we adhere to the arXiv:Web:Code ratio of 2:4:1 recommended in \cite{llemma}.
\end{itemize}

\subsubsection{Training Setting}
\label{sec:quality-policy}
We conduct mathematical domain adaptation on a broadly pre-trained language model consisting of 1.3B parameters, adhering to the same architecture utilized by DeepSeek LLMs \citep{deepseek-llm}, hereafter referred to as DeepSeek-LLM 1.3B. Distinct versions of the model are independently trained on each mathematical dataset, each for a total of 150B tokens. All training procedures leverage the resource-efficient HAI-LLM \citep{haillm} framework. Consistent with DeepSeek LLMs training protocols, we employ the AdamW optimizer \citep{adamW}, specifying hyperparameters $\beta_1=0.9$, $\beta_2=0.95$, and a $\mathrm{weight\_decay}$ of 0.1. The learning rate follows a multi-stage schedule: an initial linear warmup lasting 2,000 steps leads to the peak learning rate, which is then decayed to 31.6\% of its maximum after 80\% of total steps, and further reduced to 10.0\% by the completion of 90\% of training. The peak learning rate is fixed at 5.3e-4. Throughout, we utilize a batch size accommodating 4M tokens and set the context window to 4K tokens.

\subsubsection{Evaluation Results}

\textbf{DeepSeekMath~Corpus is distinguished by its high quality, substantial scale, and comprehensive multilingual representation of mathematical material.}

\begin{itemize}[topsep=0pt]
    \item \textbf{High-quality}:
    We assess downstream task performance across eight mathematical benchmarks using a few-shot chain-of-thought prompting approach \cite{cot}. As reported in Table \ref{tab:corpora_comparison}, models fine-tuned on DeepSeekMath~Corpus consistently outperform those trained on alternative corpora. Figure \ref{fig:corpora_comparisons} further demonstrates that a model trained with DeepSeekMath~Corpus surpasses Proof-Pile-2 at the 50B token mark (equivalent to one epoch over Proof-Pile-2), suggesting that the DeepSeekMath~Corpus provides superior average data quality.
    \item \textbf{Multilingual}:
    The DeepSeekMath~Corpus contains resources in several languages, with English and Chinese serving as the two primary languages. Data from Table \ref{tab:corpora_comparison} indicates that exposure to DeepSeekMath~Corpus enhances mathematical reasoning for both English and Chinese, whereas most available mathematical datasets—dominated by English content—exhibit limited benefits and can even impair performance for mathematical reasoning tasks in Chinese.
    \item \textbf{Large-scale}:
    In terms of volume, the DeepSeekMath~Corpus significantly exceeds prior mathematical corpora. As visualized in Figure \ref{fig:corpora_comparisons}, models such as DeepSeek-LLM 1.3B, when trained on DeepSeekMath~Corpus, demonstrate faster progression in learning and exhibit more enduring gains. By contrast, alternative corpora, being much more limited in size, require frequent repetition during training, which leads to model improvement plateauing at an earlier stage.
\end{itemize}

\subsection{Training and Evaluating DeepSeekMath-Base 7B}

This section details the DeepSeekMath-Base 7B model, a foundational system engineered for robust mathematical reasoning. The model is initialized using DeepSeek-Coder-Base-v1.5 7B \citep{deepseek-coder} and trained on 500B tokens. Data composition consists of 56\% DeepSeekMath Corpus, 4\% AlgebraicStack, 10\% arXiv, 20\% Github code, with the final 10\% drawn from Common Crawl’s English and Chinese natural language data. The majority of training settings align with Section \ref{sec:quality-policy}, with the exceptions being a peak learning rate of 4.2e-4 and a batch size set at 10M tokens.

To rigorously evaluate DeepSeekMath-Base 7B, we analyze its mathematical proficiency through multiple lenses: autonomous solution generation without auxiliary tools, mathematical problem solving leveraging external resources, and formal theorem proving capabilities. Furthermore, we outline the model’s general competencies, such as natural language comprehension, logical reasoning, and programming aptitude.

\paragraph{Stepwise Mathematical Problem Solving}

DeepSeekMath-Base’s problem-solving abilities are benchmarked using few-shot chain-of-thought prompting \citep{cot} across eight diverse English and Chinese test suites. The evaluation covers benchmarks for quantitative reasoning (such as GSM8K \citep{gsm8k}, MATH \citep{MATH}, CMATH \citep{wei2023cmath}) and multiple-choice tests (including MMLU-STEM \citep{mmlu} and Gaokao-MathQA \citep{agieval}), thus spanning a wide mathematical spectrum from elementary school to university-level challenges.

As detailed in Table \ref{tab:base_math_cot}, DeepSeekMath-Base 7B achieves top scores on all eight evaluation benchmarks when compared to other open-source base models, which include both the popular Mistral 7B \citep{mistral} and the newly introduced Llemma 34B \citep{llemma}, the latter of which was mathematically trained with Proof-Pile-2 \citep{llemma}. Of particular significance, DeepSeekMath-Base registers a more than 10\% absolute gain over previous open-source base models on the competitive MATH dataset, and even surpasses Minerva 540B \citep{minerva}, a closed-source model with a 77-fold increase in size, which was developed from PaLM \citep{palm} and subjected to additional math text training.

\paragraph{Mathematical Problem Solving with Tool Use} To assess program-assisted mathematical reasoning, we use few-shot program-of-thought prompting \citep{pot,pal} on GSM8K and MATH. Here, models are instructed to tackle each task by generating Python code that may invoke libraries such as \textit{math} or \textit{sympy} for complex operations, with the program's final execution output regarded as the predicted solution. According to Table \ref{tab:base_math_pot_proof}, DeepSeekMath-Base 7B achieves better results than the previous leading model, Llemma 34B.

\paragraph{Formal Mathematics}

Automating formal proofs helps reinforce both the validity and efficiency of mathematical argumentation, a field that has seen growing interest lately. We assess DeepSeekMath-Base 7B's ability for informal-to-formal proof conversion following the task described in \citep{dsp_proof}, where the objective is to output a formal proof given an informal conjecture, its formalized form, and an accompanying informal explanation. Evaluations are conducted on the miniF2F benchmark \citep{minif2f}, which targets formalized Olympiad-level mathematics, by prompting models to output Isabelle-formal proofs for each item using few-shot examples. Following procedures from \cite{dsp_proof}, generated proof outlines are completed by the Sledgehammer automated prover \citep{sledgehammer} to fill any missing inferential steps. As indicated in Table \ref{tab:base_math_pot_proof}, DeepSeekMath-Base 7B showcases impressive performance for proof autoformalization.

\paragraph{Natural Language Understanding, Reasoning, and Code}

We further benchmark model abilities in natural language comprehension using MMLU \citep{mmlu}, reasoning via BBH \citep{bbh}, and programming aptitude with HumanEval \citep{codex} and MBPP \citep{mbpp}. Table \ref{tab:understanding_reasoning_code} demonstrates that DeepSeekMath-Base 7B makes meaningful progress on MMLU and BBH relative to its predecessor, DeepSeek-Coder-Base-v1.5 \citep{deepseek-coder}, underlining the benefit of math-focused training for language understanding and reasoning skills. Simultaneously, incorporating code data for continual training allows DeepSeekMath-Base 7B to sustain the coding benchmark performance of DeepSeek-Coder-Base-v1.5. Collectively, DeepSeekMath-Base 7B delivers substantially better results than the general-purpose Mistral 7B \citep{mistral} across all three reasoning and programming tasks.

\section{Supervised Fine-Tuning}

\subsection{SFT Data Curation}

We develop a dataset for mathematical instruction-tuning that includes both English and Chinese questions, sourced from a range of mathematical domains and encompassing multiple levels of difficulty. Each question is associated with a corresponding solution utilizing one of three methodologies: chain-of-thought (CoT) \citep{cot}, program-of-thought (PoT) \citep{pot,pal}, or tool-integrated reasoning formats \citep{tora}. The compiled dataset contains a total of 776,000 training instances.

\begin{itemize}[topsep=0pt]
    \item \textbf{English mathematical datasets}: We enhance GSM8K and MATH datasets with tool-integrated answers, select a subset from MathInstruct \citep{MathInstruct}, and utilize the training data from Lila-OOD \citep{lila}, which features solutions formulated through CoT or PoT approaches. The English component encompasses a broad spectrum of mathematical areas, such as algebra, probability, number theory, calculus, and geometry.
    \item \textbf{Chinese mathematical datasets}: Our curation involves aggregating Chinese K-12 mathematics questions, covering 76 distinct topics including linear equations, and annotating solutions using both CoT and tool-integrated reasoning strategies.
\end{itemize}

\subsection{Training and Evaluating DeepSeekMath-Instruct 7B}

This section presents DeepSeekMath-Instruct 7B, which is fine-tuned for mathematical instruction using DeepSeekMath-Base as the foundation. During training, multiple examples are concatenated in a randomized manner up to a maximum of 4,000 tokens for each context. The training regime consists of 500 steps, a batch size of 256, and employs a fixed learning rate of 5e-5.

We assess the mathematical reasoning abilities of the models under two conditions—with and without tool utilization—across four quantitative reasoning benchmarks that include both English and Chinese. To contextualize our results, we compare our model's performance against state-of-the-art models available at the time:

\begin{itemize}[topsep=0pt]
    \item \textbf{Closed-source models} evaluated consist of: (1) the GPT family, especially the advanced GPT-4 \citep{gpt4} and GPT-4 Code Interpreter~\footnote{\url{https://openai.com/blog/chatgpt-plugins\#\#code-interpreter}}; (2) Gemini Ultra and Pro \citep{gemini}; (3) Inflection-2 \citep{inflection-2}; (4) Grok-1~\footnote{\url{https://x.ai/model-card}}; and models recently introduced by Chinese companies, such as (5) Baichuan-3~\footnote{\url{https://www.baichuan-ai.com}}; and (6) the most recent GLM-4~\footnote{\url{https://open.bigmodel.cn/dev/api\#glm-4}} from the GLM series \citep{glm}.
\end{itemize}

These models are intended for broad applicability and have typically passed through multiple alignment stages.

\item \textbf{Open-source models} comprise both general-purpose and mathematically-augmented systems, including: (1) DeepSeek-LLM-Chat 67B \citep{deepseek-llm}, (2) Qwen 72B \citep{qwen}, (3) SeaLLM-v2 7B \citep{seallm}, and (4) ChatGLM3 6B \citep{chatglm3}. The mathematically specialized models are: (5) InternLM2-Math 20B~\footnote{\url{https://github.com/InternLM/InternLM-Math}}, an extension of InternLM2 refined through math-specific training and instruction tuning; (6) Math-Shepherd-Mistral 7B, which leverages PPO training \citep{schulman2017proximal} applied to Mistral 7B \citep{mistral} using a process-supervised reward model; (7) the WizardMath series \citep{wizardmath}, designed to enhance mathematical reasoning for Mistral 7B and Llama-2 70B \citep{llama2} via evolve-instruct—a method of instruction tuning relying on AI-generated instructions—and PPO training, primarily utilizing datasets like GSM8K and MATH; (8) MetaMath 70B \citep{metamath}, which involves Llama-2 70B fine-tuned with augmented GSM8K and MATH data; (9) ToRA 34B \cite{tora}, which is CodeLlama 34B tailored for tool-integrated mathematical reasoning tasks; and (10) MAmmoTH 70B \citep{MathInstruct}, which comprises Llama-2 70B instruction-tuned on the MathInstruct dataset.
\end{itemize}

Table \ref{tab:sft_rl_math} illustrates that, when tool usage is not permitted during evaluation, DeepSeekMath-Instruct 7B achieves impressive performance in stepwise reasoning. Particularly on the challenging MATH benchmark, this model outperforms every open-source model and most proprietary systems (such as Inflection-2 and Gemini Pro) by a margin of at least 9\% in absolute accuracy. This holds even when compared to models of much greater scale (like Qwen 72B) or those further refined via math-oriented reinforcement learning techniques (such as WizardMath-v1.1 7B). Although DeepSeekMath-Instruct matches the results of leading Chinese proprietary models GLM-4 and Baichuan-3 on MATH, it still lags behind GPT-4 and Gemini Ultra.

When models are evaluated under conditions allowing both natural language reasoning and tool-assisted (i.e., programmatic) methods, DeepSeekMath-Instruct 7B attains an accuracy approaching 60\% on MATH—outperforming all other open-source contenders. On additional benchmarks, our model performs competitively with DeepSeek-LLM-Chat 67B, the prior state-of-the-art, despite being one-tenth the size.

\section{Reinforcement Learning}

\subsection{Group Relative Policy Optimization} It has been well established that reinforcement learning (RL) can further refine the mathematical reasoning capabilities of LLMs after the Supervised Fine-Tuning (SFT) phase \citep{wang2023math,wizardmath}. In the following, we present our novel and efficient RL approach: Group Relative Policy Optimization (GRPO).

\subsubsection{From PPO to GRPO} Proximal Policy Optimization (PPO) \citep{schulman2017proximal} is a leading actor-critic RL method, frequently employed in the RL fine-tuning of LLMs \citep{ouyang2022training}. PPO updates LLMs by maximizing the surrogate objective defined as:

\begin{equation}
\footnotesize
    \mathcal{J}_{PPO}(\theta) = \mathbb{E}{[q \sim P(Q), o \sim \pi_{\theta_{old}}(O|q)]} \frac{1}{|o|} \sum_{t=1}^{|o|} \min \left[ \frac{\pi_\theta(o_{t} | q, o_{<t})}{\pi_{\theta_{old}}(o_{t} | q, o_{<t})} A_{t}, \text{clip} \left( \frac{\pi_\theta(o_{t} | q, o_{<t})}{\pi_{\theta_{old}}(o_{t} | q, o_{<t})}, 1 - \varepsilon,

In this context, $\pi_{\theta}$ and $\pi_{\theta_{old}}$ refer to the current and previous policy models, respectively, while $q$ and $o$ denote questions and outputs that are drawn from the question dataset and generated by the prior policy $\pi_{\theta_{old}}$. The parameter $\varepsilon$ serves as a clipping hyper-parameter that is introduced by PPO to improve training stability. The advantage term $A_t$ is obtained via the Generalized Advantage Estimation (GAE) algorithm \citep{gae}, which leverages the observed rewards $\{r_{\ge t}\}$ and predictions from a learned value function $V_{\psi}$. Consequently, PPO necessitates learning a value function in addition to training the policy model. To curb excessive optimization of the reward model, a commonly employed method is to include a token-level KL penalty, measured with respect to a reference model, in the reward at every step \citep{ouyang2022training}, specifically:

\begin{equation}
    r_{t} = r_\varphi(q, o_{\le t}) - \beta \log\frac{\pi_{\theta}(o_{t}|q, o_{<t})}{\pi_{ref}(o_{t}|q, o_{<t})},
\label{eq:PPO-reward}
\end{equation}

Here, $r_\varphi$ designates the reward model, $\pi_{ref}$ stands for the reference model (typically the initial SFT model), and $\beta$ is the hyper-parameter controlling the strength of the KL penalty.

A notable challenge is that the value function in PPO generally requires a model of similar scale to the policy network, which incurs significant memory and computation costs. Further, the value function is primarily used during RL training to reduce the variance in the computed advantages. In LLM-based scenarios, it is typical that only the final token receives a reward from the reward model, complicating efforts to construct a value function with reliable predictions at every token position. To overcome these challenges, we introduce Group Relative Policy Optimization (GRPO), as illustrated in Figure \ref{fig:grpo}. GRPO removes the necessity for an extra value function by instead taking the mean reward of several output samples, all prompted by the same question, to serve as the baseline. Concretely, for each question $q$, GRPO draws a set of outputs $\{o_1, o_2, \cdots, o_G\}$ from the old policy $\pi_{\theta_{old}}$ and then updates the policy network by maximizing the following criterion:

\begin{equation}
\footnotesize
\begin{split}
    \mathcal{J}_{GRPO}(\theta) &= \mathbb{E}{[q \sim P(Q), \{o_i\}_{i=1}^G \sim \pi_{\theta_{old}}(O|q)]}  \\
    & \frac{1}{G}\sum_{i=1}^G\frac{1}{|o_i|} \sum_{t=1}^{|o_i|} \left\{ \min \left[ \frac{\pi_\theta(o_{i,t} | q, o_{i,<t})}{\pi_{\theta_{old}}(o_{i,t} | q, o_{i,<t})} \hat{A}_{i,t}, \text{clip} \left( \frac{\pi_\theta(o_{i,t} | q, o_{i,<t})}{\pi_{\theta_{old}}(o_{i,t} | q, o_{i,<t})}, 1 - \varepsilon, 1 + \varepsilon \right)  \hat{A}_{i,t} \right] - \beta \mathbb{D}_{KL}\left[\pi_{\theta} || \pi_{ref}\right]\right\} ,
\end{split}
\label{eq:GRPO-obj}
\end{equation}

In this objective, $\varepsilon$ and $\beta$ denote hyper-parameters, while $\hat{A}_{i,t}$ indicates the advantage score determined by examining the relative rewards within the output group, with specifics provided in subsequent subsections. This group-relative mechanism for estimating advantages complements the reward model’s preference-based framework, as reward models are often learned from datasets containing comparisons of multiple outputs for each question. Furthermore, unlike PPO’s KL penalty term (added inside the reward), GRPO enforces KL regularization by incorporating the KL divergence between the current policy and the reference policy directly in the loss term, thus keeping the computation of $\hat{A}_{i,t}$ more straightforward. For estimating the KL divergence (as opposed to the penalty in (\ref{eq:PPO-re

\begin{algorithm}[t]
  \small
  \caption{Iterative Group Relative Policy Optimization}
  \textbf{Input} initial policy model $\pi_{\theta_{\text{init}}}$; reward models $r_\varphi$; task prompts $\mathcal{D}$; 
  hyperparameters $\varepsilon$, $\beta$, $\mu$
  \begin{algorithmic}[1]
    \State Initialize policy model $\pi_\theta$ with $\pi_{\theta_{\text{init}}}$
    \For{iteration = 1, \dots, I}
      \State Assign reference model $\pi_{ref} \gets \pi_{\theta}$
      \For{step = 1, \dots, M}
        \State Draw a batch $\mathcal{D}_b$ from dataset $\mathcal{D}$
        \State Set the previous policy model $\pi_{\theta_{old}} \gets \pi_{\theta}$ 
        \State For every prompt $q$ in $\mathcal{D}_b$, generate $G$ responses $\{o_i\}_{i=1}^G \sim \pi_{\theta_{old}} (\cdot \mid q)$
        \State Evaluate each generated response $o_i$ using $r_\varphi$ to produce rewards $\{r_i\}_{i=1}^{G}$
        \State For the $t$-th token in $o_i$, determine $\hat{A}_{i,t}$ using group relative advantage calculation
        \For{GRPO iteration = 1, \dots, $\mu$}
          \State Adjust the policy model $\pi_\theta$ to maximize the GRPO objective (see Equation \ref{eq:GC-GRPO})
        \EndFor
      \EndFor 
      \State Continuously train $r_\varphi$ with a replay buffer.
    \EndFor 
  \end{algorithmic}
  \textbf{Output} $\pi_\theta$
  \label{alg:iter-grpo}
\end{algorithm}

\subsubsection{Outcome Supervision RL with GRPO} Consider a question $q$ where a set of candidate outputs $\{o_1, o_2, \cdots, o_G\}$ is generated from the former policy $\pi_{\theta_{old}}$. Each output is evaluated using the reward model to obtain $G$ reward values, denoted as $\mathbf{r}=\{r_1, r_2, \cdots, r_G\}$. These raw rewards undergo normalization, wherein the mean of the group is subtracted and the result is divided by the group’s standard deviation. Under outcome supervision, each output $o_i$ receives the normalized reward upon completion; all tokens within $o_i$ receive this same normalized score as their advantage, mathematically: $\hat{A}_{i, t} = \widetilde{r}_i = \frac{r_i- {\rm mean}(\mathbf{r})}{{\rm std}(\mathbf{r})}$. The policy parameters are updated by maximizing the objective in Equation (\ref{eq:GRPO-obj}).

\subsubsection{Process Supervision RL with GRPO} In outcome supervision, the reward is limited to the end of each output, which may not offer sufficiently granular feedback for complex reasoning or mathematical problems. As in \cite{wang2023math}, process supervision is adopted, granting a reward at each intermediate reasoning step. Specifically, given question $q$ and a set of $G$ sampled outputs $\{o_1, o_2, \cdots, o_G\}$, a process reward model scores each step within these outputs, yielding reward collections $\mathbf{R} = \{ \{r_1^{index(1)},\cdots,r_1^{index(K_1)}\}, \cdots,  \{r_G^{index(1)},\cdots,r_G^{index(K_G)}\} \}$, with $index(j)$ marking the terminal token of step $j$, and $K_i$ representing the count of steps in $o_i$. These stepwise rewards are normalized similarly: $\widetilde{r}_i^{index(j)} = \frac{r_i^{index(j)} - {\rm mean(\mathbf{R})}}{{\rm std(\mathbf{R})}}$. Then, for every token at position $t$ in $o_i$, the advantage is defined as the sum of all subsequent normalized step rewards: $\hat{A}_{i, t} = \sum_{index(j) \ge t} \widetilde{r}_i^{index(j)}$. Optimization of the policy is subsequently driven by maximizing the objective in Equation (\ref{eq:GR

\subsubsection{Iterative RL with GRPO} As reinforcement learning advances, the previously trained reward model may become inadequate for supervising the evolving policy model. To address this, we examine an iterative RL approach employing GRPO. As depicted in Algorithm \ref{alg:iter-grpo}, iterative GRPO involves generating fresh training datasets for the reward model by utilizing samples produced by the policy model. We continue to update the reward model through a replay strategy that includes 10\% of data from prior iterations. Subsequently, the policy model serves as the reference for further updates, and we iteratively refine the policy model under the guidance of the updated reward model.

\subsection{Training and Evaluating DeepSeekMath-RL}

Our RL training is built on DeepSeekMath-Instruct 7B. The dataset for RL consists of chain-of-thought questions drawn from the SFT portion of GSM8K and MATH, totaling approximately 144K items. We omit other SFT questions to assess how RL training influences benchmark performance when those datasets are not represented during RL. The reward model's training data is prepared in accordance with \citep{wang2023math}. The initial reward model is trained from DeepSeekMath-Base 7B with a learning rate of 2e-5. For GRPO, the policy model's learning rate is set at 1e-6 and the KL penalty is 0.04. Each question is used to generate $64$ candidate responses. The sequence length maximum is set to 1024, with a batch size also at 1024. The policy model undergoes one update following each exploration cycle. DeepSeekMath-RL 7B is evaluated on benchmark datasets in the same way as DeepSeekMath-Instruct 7B. In the case of DeepSeekMath-RL 7B, benchmarks featuring GSM8K and MATH with chain-of-thought reasoning are classified as in-domain, while all remaining benchmarks are treated as out-of-domain.

Table \ref{tab:sft_rl_math} presents a comparison of open- and closed-source models on English and Chinese evaluation tasks, assessing both chain-of-thought and tool-assisted reasoning approaches. Key observations include: 1) When applying chain-of-thought reasoning, DeepSeekMath-RL 7B achieves 88.2\% accuracy on GSM8K and 51.7\% on MATH. These results exceed those of all open-source models within the 7B–70B parameter range, and also surpass most closed-source alternatives. 2) Notably, the training for DeepSeekMath-RL 7B is limited strictly to chain-of-thought instruction tuning data derived from GSM8K and MATH, building directly on DeepSeekMath-Instruct 7B. Even with this narrow training dataset, DeepSeekMath-RL 7B demonstrates improved performance across all assessment benchmarks when compared to DeepSeekMath-Instruct 7B, highlighting the efficacy of reinforcement learning.

\section{Discussion}
This section summarizes insights gained from our work in pre-training and reinforcement learning (RL) experiments.

\subsection{Lessons Learnt in Pre-Training}

We begin by recounting experiences from the pre-training process. Except where indicated otherwise, all experiments follow the training protocols established in Section \ref{sec:quality-policy}. In the context of this discussion, the term DeepSeekMath Corpus refers to an 89B-token collection assembled during the second data gathering cycle.

\subsubsection{Code Training Benefits Mathematical Reasoning}

It has been popularly conjectured, but not empirically proven, that exposure to code during training enhances reasoning abilities. Here, we provide evidence—specifically within mathematics—that code training bolsters a model's mathematical reasoning skills, applicable both with and without tool integration.

To examine the impact of code training on mathematical reasoning, we designed experiments employing both a two-stage and a one-stage training paradigm:

\textbf{Two-Stage Training}

\begin{itemize}[topsep=0pt]
    \item \textbf{Code Training for 400B Tokens $\rightarrow$ Math Training for 150B Tokens}: The DeepSeek-LLM 1.3B model undergoes initial training on 400B code tokens, followed by an additional 150B tokens consisting of mathematics-related data.
    \item \textbf{General Training for 400B Tokens $\rightarrow$ Math Training for 150B Tokens}: For comparison, we conduct a baseline experiment where the first training phase utilizes 400B general-purpose tokens (drawn from a comprehensive general corpus prepared by DeepSeek-AI) in lieu of code tokens. This aims to evaluate whether code-specific pretraining confers advantages over general corpus pretraining in mathematical reasoning tasks.
\end{itemize}

\textbf{One-Stage Training}

\begin{itemize}[topsep=0pt]
    \item \textbf{Math Training for 150B Tokens}: Here, DeepSeek-LLM 1.3B is trained exclusively on 150B tokens from mathematical datasets.
    \item \textbf{Training on a mixture of 400B Code Tokens and 150B Math Tokens}: Sequential math training after code pretraining has been observed to diminish coding capabilities. Thus, we investigate an alternative approach: jointly training the model on a combined dataset of code and math tokens, with the dual objective of fostering mathematical reasoning and alleviating catastrophic forgetting.
\end{itemize}

\paragraph{Results}
The outcomes under varied training regimes are presented in Table \ref{tab:code-math} and Table \ref{tab:code-math-general-eval}.

Introducing code token pretraining proves beneficial for program-aided mathematical reasoning in both two-stage and integrated training setups. Table \ref{tab:code-math} reveals that within the two-stage protocol, code pretraining alone provides substantial gains in the model’s ability to address GSM8K and MATH tasks via Python. Subsequent training on mathematical data yields further improvements. In the one-stage setting, co-training on mixed code and math tokens effectively addresses catastrophic forgetting—a shortcoming observed with purely sequential training—and simultaneously enhances both coding performance (see Table \ref{tab:code-math-general-eval}) and mathematical reasoning supported by programming (Table \ref{tab:code-math}).

Beyond program-aided tasks, code token pretraining also supports mathematical reasoning in the absence of external tools. In two-stage training, code-focused pretraining alone yields moderate performance gains and primes the model for more efficient subsequent math training, ultimately leading to superior results. Conversely, simultaneous exposure to code and math tokens in a single-stage regime appears to hinder mathematical reasoning when tool use is not permitted. We posit that the 1.3B parameter size of DeepSeek-LLM may be insufficient for the model to effectively internalize and synthesize both code and mathematical content at once.

\subsubsection{ArXiv Papers Appear to Offer Limited Benefits for Mathematical Reasoning}

ArXiv papers are frequently incorporated as a core element in the construction of mathematical pre-training datasets \citep{minerva,gpt-f,llemma,mathpile}. Yet, a comprehensive evaluation of their specific contribution to mathematical reasoning performance has been largely absent in prior research. Surprisingly, our experimental findings indicate that pre-training on arXiv content alone does not enhance mathematical reasoning. Our investigation encompasses a range of model architectures, including DeepSeek-LLM 1.3B and DeepSeek-Coder-Base-v1.5 7B \citep{deepseek-coder}, leveraging arXiv data processed via two distinct methodologies:

\begin{itemize}[topsep=0pt]
    \item \textbf{MathPile} \citep{mathpile}: a dataset with 8.9B tokens, filtered and cleaned according to heuristic rules, with more than 85\% consisting of scientific arXiv papers;
    \item \textbf{ArXiv-RedPajama} \citep{redpajama}: a 28.0B-token collection comprising arXiv LaTeX files after the removal of preambles, comments, macros, and bibliographic entries.
\end{itemize}

We conduct separate pre-training runs with DeepSeek-LLM 1.3B (for 150B tokens) and DeepSeek-Coder-Base-v1.5 7B (for 40B tokens) exclusively on each arXiv-based dataset. Results consistently demonstrate that relying on arXiv papers for pre-training fails to yield measurable gains—and occasionally even harms performance—on several mathematical evaluation suites of varying difficulty included in our analysis. These evaluations span quantitative reasoning tasks (e.g., GSM8K and MATH, see Table \ref{tab:arxiv-cot}), multiple-choice STEM challenges (e.g., MMLU-STEM, Table \ref{tab:arxiv-cot}), and formal mathematics (e.g., miniF2F, Table \ref{tab:arxiv_atp}).

Nonetheless, this outcome should be interpreted with caution. Several relevant questions remain unexamined, including:

\begin{itemize}[topsep=0pt]
    \item How arXiv-derived tokens may affect particular math domains not assessed here, such as converting formal mathematical content to informal exposition (informalization);
    \item The influence of arXiv data when integrated with other types of training corpora;
    \item Whether scaling up model size could reveal advantages of training on arXiv papers.
\end{itemize}

Consequently, more detailed studies are needed, which we defer to subsequent research.

\subsection{Insights of Reinforcement Learning}
\subsubsection{Advancing Towards a Unified Paradigm}
Here, we propose a comprehensive framework to study and unify several training strategies—including SFT, RFT, DPO, PPO, and GRPO—and experimentally probe various factors within this paradigm. Generally, the parameter gradient for any given training method can be formulated as follows:

\begin{equation}
    \nabla_{\theta}\mathcal{J}_{\textcolor{red}{\mathcal{A}}}(\theta) = \mathbb{E}[\underbrace{(q,o) \sim \textcolor{red}{\mathcal{D}}}_{Data \ Source}]\left( \frac{1}{|o|} \sum_{t=1}^{|o|}  \underbrace{GC_{{\mathcal{A}}}(q, o, t, \textcolor{red}{\pi_{{rf}}})}_{Gradient \ Coefficient}  \nabla_{\theta}\log \pi_{\theta}(o_t | q, o_{<t})\right).
\label{eq:objective}
\end{equation}

This formulation relies on three central elements: 1) \textit{Data Source $\mathcal{D}$}, which specifies the origin of training examples; 2) \textit{Reward Function $\pi_{{rf}}$}, responsible for providing the feedback signal used in training; and 3) \textit{Algorithm $\mathcal{A}$}, which takes in both the training data and the reward signals and produces the gradient coefficient $GC$ used to determine the extent of adjustment—either penalization or reinforcement—for the data. We discuss several prominent approaches following this standardized framework:

\begin{itemize}
    \item  \textbf{Supervised Fine-tuning (SFT)}: This approach adapts a pretrained model to human-annotated SFT datasets by additional supervised training.
    \item \textbf{Rejection Sampling Fine-tuning (RFT)}: RFT involves further fine-tuning the model that underwent SFT, utilizing filtered response samples generated from the SFT model, retaining only outputs deemed correct with respect to SFT prompts.
    \item \textbf{Direct Preference Optimization (DPO)}: DPO improves upon the SFT model by conducting additional fine-tuning on expanded response data sampled from the SFT model, leveraging a pairwise DPO loss.
    \item \textbf{Online Rejection Sampling Fine-tuning (Online RFT)}: Unlike conventional RFT, Online RFT starts from the SFT-initialized policy and updates it through further fine-tuning, drawing augmented outputs sampled from the ongoing, real-time policy model.
    \item \textbf{PPO/GRPO}: With these methods, the policy model is initialized using the SFT model and is subsequently reinforced through interaction-based training, utilizing responses produced by the real-time policy.
\end{itemize}

The elements constituting these approaches are outlined in Table \ref{tab:data_policy}. For a comprehensive exposition of the derivation procedures, please consult Appendix \ref{app:analysis-rl}.

\paragraph{Observation about Data Source} We categorize the data sources into two primary types: online and offline sampling. Online sampling indicates that the training set is generated from the ongoing explorations conducted by the policy model in real time, whereas offline sampling refers to data collected through the sampling of the initial SFT model. RFT and DPO utilize the offline approach, while Online RFT and GRPO adopt the online sampling methodology.

From Figure \ref{fig:rl-analysis}, it is evident that Online RFT exhibits markedly better performance than RFT across both evaluation benchmarks. Initially, the effectiveness of Online RFT mirrors that of RFT, but as training progresses, Online RFT achieves a clear advantage, underscoring the benefits of online data acquisition. This phenomenon can be explained by noting that, at the start, the actor and the SFT model behave similarly, so the sampled outputs do not substantially differ. However, as training advances, the actor-generated samples start to diverge more noticeably, at which point access to up-to-date data provides a substantial edge.

\paragraph{Observation about Gradient Coefficient} The algorithms convert the reward signals into gradient coefficients for updating model parameters. In our experimental setup, the reward function is classified as either `Rule' or `Model.' The `Rule' category pertains to determining the correctness of a response strictly by whether the answer is right or wrong. The `Model' variant involves employing a reward model trained on rule-based evaluations to assign a score to each response. Equations \ref{eq:GC-RFT} and \ref{eq:GC-GRPO} delineate a fundamental distinction between GRPO and Online RFT: GRPO is characterized by a gradient coefficient that is directly modulated by the output of the reward model, allowing it to assign variable reinforcement or penalties depending on the magnitude of the reward. Conversely, Online RFT does not feature this mechanism; it applies equal reinforcement to all correct responses, and does not penalize incorrect responses.

As illustrated in Figure \ref{fig:rl-analysis}, GRPO demonstrates superior performance compared to online RFT, underscoring the effectiveness of modifying the coefficients associated with positive and negative gradients. Moreover, the results show that GRPO+PS achieves better outcomes than GRPO+OS, emphasizing the advantages of applying step-aware, fine-grained gradient coefficients. Additionally, we investigate the impact of iterative RL by performing two successive iterations in our experiments. Figure \ref{fig:iter-rl} reveals that iterative RL leads to a marked performance gain, most notably during the initial iteration.

\subsubsection{Why RL Works?} In this work, reinforcement learning is conducted using a subset derived from instruction tuning datasets, resulting in marked improvements over the instruction-tuned baseline. To elucidate the effectiveness of reinforcement learning, we analyze Pass@K and Maj@K accuracies for both Instruct and RL-optimized models on two evaluation sets. Figure \ref{fig:maj-pass} demonstrates that while RL increases Maj@K accuracy, it does not enhance Pass@K. This pattern suggests that reinforcement learning makes the output distribution more robust, enhancing the likelihood that the correct answer appears within the TopK responses, but does not necessarily improve the model’s core skills. In essence, \textbf{the main gain arises from an increased probability of correct answers appearing in TopK, rather than the improvement of underlying abilities.} Similarly, \citep{wang2023making} report a \textbf{misalignment problem} in reasoning performance with SFT models, and show that introducing various preference alignment techniques \citep{yuan2023rrhf,song2023preference,wang2023making} leads to better reasoning capabilities in these systems.

\subsubsection{How to Achieve More Effective RL?}
\label{sec:effec_rl}
Our results establish the effectiveness of reinforcement learning for mathematical reasoning tasks and offer a unifying framework to interpret leading training methods. Within this framework, these approaches can all be seen as forms of direct or simplified RL. As highlighted in Equation \ref{eq:objective}, three fundamental aspects shape these methods: Data Source, Algorithm, and Reward Function. We discuss prospective advancements related to each component.

\paragraph{Data Source} The selection of data is foundational for all training paradigms. For RL specifically, the data source refers to unlabeled question prompts coupled with responses generated from the policy model. In our study, the data was limited to prompts from instruction tuning and responses obtained via basic nucleus sampling. We believe this may explain why RL in our case only benefited Maj@K. Looking forward, we intend to expand our RL framework to handle out-of-distribution question prompts and incorporate \textbf{more sophisticated sampling (decoding) methods}, including those utilizing tree-search algorithms \citep{yao2023tree}. Furthermore, the adoption of \textbf{efficient inference techniques} \citep{xia-etal-2023-speculative,leviathan2023fast,kwon2023efficient,xia2024unlocking}—critical for improving the exploration efficiency of policy models—will be a focal point of future work.

\paragraph{Algorithms} Algorithms leverage both data and reward signals to compute gradient coefficients, which are then used to adjust the parameters of the model. According to Equation \ref{eq:objective}, contemporary approaches, to varying degrees, fundamentally \textbf{TRUST} the feedback provided by the reward function to modulate the conditional probability of producing specific tokens. Nevertheless, the infallibility of the reward signal cannot be assumed, particularly in highly intricate scenarios. For instance, the PRM800K datasets \citep{lightman2023let}, despite being meticulously labeled by proficient annotators, still exhibit an error rate of about 20\% in their annotations\footnote{\url{https://github.com/openai/prm800k/issues/12\#issuecomment-1728491852}}. Therefore, we propose to investigate reinforcement learning algorithms capable of maintaining robustness in the face of noisy or unreliable reward feedback. Our perspective is that such \textbf{WEAK-TO-STRONG} \citep{burns2023weak} alignment strategies could profoundly transform the design of learning algorithms.

\paragraph{Reward Function} The reward function serves as the origin of learning signals during training. In reinforcement learning, the reward is commonly delivered through a neural reward model. We identify three pivotal avenues for advancing reward models: 1) \textbf{Improving the reward model’s generalization.} The reward model should be capable of effectively transferring to out-of-distribution tasks and complex generation outputs; otherwise, reinforcement learning may merely reinforce current LLM behaviors without catalyzing substantive performance gains; 2) \textbf{Capturing the uncertainty of reward models.} Incorporating uncertainty estimates may serve as an essential interface connecting imperfect reward models with weak-to-strong learning frameworks; 3) \textbf{Constructing high-quality, process-level reward models efficiently} that are able to supply detailed feedback throughout multi-step reasoning tasks \citep{lightman2023let,wang2023math}.

\section{Conclusion, Limitation, and Future Work} In this paper, we introduce DeepSeekMath, which surpasses all open-source models on the challenging MATH benchmark and delivers performance nearing that of proprietary models. DeepSeekMath~is based on DeepSeek-Coder-v1.5 7B, with ongoing training on 500B tokens—of which a substantial 120B are mathematical tokens collected from Common Crawl. Our comprehensive ablation analysis highlights that web pages are a particularly fruitful resource for acquiring high-quality mathematical content, while the arXiv dataset appears less advantageous than anticipated. Additionally, we propose Group Relative Policy Optimization (GRPO), a modification of Proximal Policy Optimization (PPO), capable of enhancing mathematical reasoning with reduced memory requirements. Empirical results confirm that GRPO remains effective, even when DeepSeekMath-Instruct 7B already performs strongly on standard benchmarks. Furthermore, we offer a cohesive framework for interpreting a range of methods and identify multiple promising avenues to advance reinforcement learning methodologies.

Despite DeepSeekMath's~strong results on quantitative reasoning evaluations, it demonstrates comparatively lower proficiency in handling geometry and theorem-proving tasks when contrasted with closed-source alternatives. For example, our preliminary testing revealed that the model struggles with questions involving triangles and ellipses, which may point to biases stemming from pre-training and fine-tuning data selection. Moreover, due to its current parameter scale, DeepSeekMath~lags behind GPT-4 in few-shot learning scenarios: while GPT-4 benefits from few-shot prompts, DeepSeekMath~exhibits comparable outcomes in both zero-shot and few-shot settings. Moving forward, we plan to refine our engineered data selection process to assemble an even higher quality pre-training dataset. We also intend to investigate the research paths outlined in Section~\ref{sec:effec_rl} to further enhance reinforcement learning techniques for large language models.

\end{document}